\cotefile{157-4}
\scanpage{196}
\mspage{70}
\deleted{\unsure{Cela} \unsure{résu}} On a bien \unsure{sur|sûr} $\Longleftarrow$, donc il reste à prouver $\Longrightarrow$. Si l'ens \unsure{d'indices} $A$ est fini, cela résulte du lemme 5.5, compte tenu que la finalité \unsure{des} $M_\alpha$ implique leur $\unsure{\pi|R|K}$-cofinalité. On doit donc simplement améliorer 5.5, par le

\uline{Lemme 5.15} \unsure{Soient} \inserted{$\pi$ un cardinal,} $I$ un ens. ordonné filtrant $\pi$-adapté, et soit $(I_\alpha)_{\alpha\in A}$ une famille de parties $I_\alpha$ $\pi$-cofinales (cf 5.8.), telle que $\operatorname{card} A\leq\pi$. Alors $\bigcap I_\alpha$ est $\pi$-cofinale.

Il suffit de prouver qu'elle est cofinale, car la propriété de stabilité par bornes supérieures de parties filtrantes de cardinal $\leq\pi$ est évidente. Quitte à remplacer $I$ par $I/i_0$ et les $I_\alpha$ par leurs traces sur $I/i_0$, la question est simplement de prouver que $\bigcap_\alpha I_\alpha\ne\varnothing$. On va construire un élément $s$ de $\bigcap I_\alpha$ par induction transfinie. Prenons un bon ordre sur $A$, et \deleted{\unsure{construisons} \unsure{une suite|un syst.}} \inserted{\unsure{prouvons} par réc. transfinie \deleted{\unsure{transfinie}}} \deleted{\uline{\unsure{croissante|croissant}} d'éléments $\unsure{i_\alpha|i'_\alpha}$ de $I$, \unsure{telle} que}
% archive sheet number: 195
\scanpage{197}
\mspage{71}
\[
 \deleted{i_\alpha\in\bigcap_{\beta\leq\alpha}I_\beta\quad\unsure{\text{pour}}\ .}
\]
\deleted{On choisit $i_0$ ($0$ élément initial de $A$) dans $I_0$ quelconque, et on}

que pour tout $\alpha\in A$, l'ensemble
\[
 J_\alpha=\bigcap_{\beta<\alpha}I_\beta
\]
est cofinal. Si on sait que $J_\alpha$ est cofinal, il en est de $=$ de $J_{\alpha+1}=J_\alpha\cap I_\alpha$, par 5.5. Donc on est ramené à prouver que $J_\alpha$ est cofinal si $\alpha$ est un ordinal limite, donc à prouver le lemme 5.15. pour une suite \inserted{transfinie} \uline{décroissante} de parties \deleted{\unsure{dans un}} $\unsure{\pi|R|K}$-cofinales, \inserted{$J_\beta$, $(\beta<\alpha)$ \unsure{quand} $\operatorname{card}\alpha\leq\pi$.} \deleted{\unsure{de}} \inserted{telle que $J_\beta=\bigcap_{\beta'<\beta}J_{\beta'}$, \deleted{\unsure{pour}} \unsure{quand} $\beta$ est limite.} \inserted{\unsure{A} prouver} que l'intersection est $\pi$-cofinale --- et on est réduit à prouver qu'elle est non vide. On va construire \unsure{par} réc. transfinie une suite \uline{croissante} d'éléments
\[
 i_\beta\in J_\beta\qquad(\beta\leq\alpha)
\]
Si $i_\beta$ est construit, comme $J_{\beta+1}$ est cofinal, il existe $i_{\beta+1}\geq i_\beta$ dans $J_{\beta+1}$, OK. Si $\beta$ \unsure{est} un ordinal limite, on prend
\[
 i_\beta=\mathop{\mathrm{b.\!sup}}\limits_{\beta'<\beta} i_{\beta'}, 
\]
$\unsure{i'_\beta|i_{\beta'}}$ est dans \deleted{\unsure{les}} chaque $I_{\beta_0}$ pour $\beta_0<\beta$,
% archive sheet number: 196
\scanpage{198}
\mspage{72}
puisque les $i_{\beta'}$ pour $\beta_0\leq\beta'<\beta$ sont \inserted{tous} dans $J_{\beta_0}\supset J_{\beta'}$, donc aussi leur limite, \deleted{\unsure{donc}} \deleted{$i_\beta$ \unsure{est}} donc
\[
 i_\beta\in\bigcap_{\beta_0<\beta}J_{\beta_0}=J_\beta .
\]
Bon la récurrence transfinie marche, et nous donne donc un \unsure{élmt|élément} \deleted{de} \unsure{$s_\alpha$} de $J_\alpha$ !

Je vais prouver aussi la \unsure{forme|formule} $\unsure{\cup|\mathrm U}$ \unsure{assertion} pour $\bigcup_{\alpha\in A}\deleted{I}M_\alpha$, quand $A$ est fini. Posant cette fois $M=\bigcup_\alpha M_\alpha$, il faut prouver, moyennant les relations $(\ast)$\inserted{\unsure{(p. 68)}}, que
\[
 \begin{aligned}
 x\in M &\Longleftrightarrow
  \text{l'ens }I_M\ (= \bigcup_\alpha I_{M_\alpha})\\
 \substack{\text{(i.e. }\exists\alpha\text{ avec}\\ x\in M_\alpha\text{)}} &\hspace{2em}\text{est cofinal dans }I
 \end{aligned}
 \tag*{$(\ast\ast)^{\sim}$}
\]
et uniquement \unsure{car|en} le premier terme de cette \unsure{équivalence} : il \unsure{doit|devrait} signifier que \deleted{$\unsure{\forall}$} $\exists\alpha$ avec $I_{M_\alpha}$ cofinal dans $I$. Bien sûr, si un des $I_{M_\alpha}$ est cofinal dans $I$, \deleted{\unsure{sa}} \inserted{\unsure{à|a} fortiori} leur réunion. Mais inversement, si la réunion est cofinale, elle est filtrante et réunion des parties (filtrantes) $I_{M_\alpha}$, dont l'une doit \unsure{bien|donc} être cofinale dans elle par le lemme 5.14, donc aussi dans $I$, qed.
% archive sheet number: 197
\scanpage{199}
\mspage{73}
\uline{Corollaire} \uline{5.16} (de prop 5.13.) Toute réunion finie de parties L-accessibles de $M$ est L-accessible. \unsure{Toute} intersection $\bigcap_{\alpha\in A}M_\alpha$ de \inserted{une famille \unsure{petite} de} parties L-accessibles de $M$, est L-accessible.

6. \unsure{Essay|Essai} avorté pour un critère d'engendrement des parties L-accessibles de $\unsure{\mathrm{Fl}|F\ell}\,M$.

Soit $M$ une catégorie munie d'une filtration cardinale \inserted{$(\mathrm{Filt}^{\unsure{\pi|R}})_{\pi\geq\pi_0}$}, et \deleted{\unsure{considérons}} soit $S$ un objet de $\mathrm{Cat}$, considérons
\[
 M(S)=\underline{\mathrm{Hom}}\,(S,M),
\]
et munissons $M(S)$ de la filtration par les
\[
 M^\pi(S)=\underline{\mathrm{Hom}}\,\bigl(S,M^\pi=\underline{\mathrm{Filt}}^\pi(M)\bigr).
\]
Je dis que c'est là une filtration cardinale de $M(S)$. \deleted{\unsure{Il}} \unsure{faut} vérifier les conditions a) b) c) de SGA 4 I 9.12. Les conditions a) ($\mathrm{Filt}^\pi$ \unsure{ess. petite|ens. petits.}) et b (condition $\mathrm{L}_{\pi_0}$, et stabilité \unsure{de} $\mathrm{Filt}^\pi$ par $\underset I{\underline{\lim}}$, \deleted{\unsure{filtran}} $I$ filtrant grand devant $\pi_0$ et de cardinal $\leq\pi$) passent aussitôt de $M$ à $M(S)$. La condition c) est plus délicate. (\unsure{représentation}
% archive sheet number: 198
\scanpage{200}
\mspage{74}
d'un objet comme $\underline{\lim}$ filtrante grande devant $\unsure{\pi|R}$ d'objets de $\mathrm{Filt}^\pi$). Je ne vais pas refaire le travail, fait \deleted{dans une \unsure{toute} \unsure{bonn}} \inserted{dans un cas} beaucoup plus général et plus délicat dans SGA 4 I §9, savoir dans le cor. 9.24 du th. \unsure{9.23|9.27|9.22}, (p. 177). La seule restriction, c'est qu'on \inserted{y} suppose les objets de $M$ \deleted{$\unsure{\pi|R}$-}accessibles, et désignant par $\pi_1$ \deleted{un} le plus petit cardinal tel que $\pi_1\geq\pi_0$ et que les objets de $\mathrm{Filt}^{\pi_0}$ soient $\pi_1$-accessibles, on se borne \inserted{\unsure{pour} la filtration cardinale de $\underline{\mathrm{Hom}}\,(S,M)$} aux cardinaux $\boxed{\pi\geq c=2^{c_0},}$ où
\[
 c_0=\mathrm{Sup}\,(\pi_1,\operatorname{Card}\mathrm{Fl}\, S).
\]
C'est \unsure{une} restriction \unsure{raisonnable} et \unsure{modérée}.

Pour le \unsure{'instant|l'instant}, je suis intéressé surtout au cas où $S=\Delta^1$, donc :
\[
 \underline{\mathrm{Hom}}\,(\Delta^1,M)=\underline{\mathrm{Fl}}\,(M),
\]
auquel cas on a $c_0=\pi_1$, donc il suffit de se borner aux cardinaux $\pi\geq 2^{\pi_1}$ pour avoir une bonne filtration cardinale sur $\underline{\mathrm{Fl}}\,(M)$.

C'est donc \unsure{l'ens|l'une} des deux conditions b), c) sur $M$ (p 44) qui est stable par passage aux $\underline{\mathrm{Hom}}\,(S,M)=M(S)$. Il en est de $=$ de la condition a) (stabilité pour petites limites ind \deleted{\unsure{filtrantes}}), et \inserted{aussi} de la condi%
% archive sheet number: 199
\scanpage{201}%
\mspage{75}%
\deleted{\unsure{tion}} d) (p. 45). De $=$ la condition d'exactitude e) (p. 46) sur les $\underrightarrow{\lim}$ filtrantes (\unsure{dérivant|devient} \unsure{d'un|donc} \unsure{lemme|jeu} \unsure{connu} \unsure{en moins|au moins}).

Pour être tranquille, je vais faire sur $M$ les hypothèses a) à e) inclus. (bien que e) me semble un peu artificielle, on devrait pouvoir l'éliminer, dans un second souffle !) Ces conditions ont donc les propriétés de stabilité souhaitables par passage aux $M(S)$, et en particulier par passage à $\mathrm{Fl}(M)$. Quitte à augmenter $\unsure{\pi|R}_0$, on peut donc supposer que les
\[
 \mathrm{Fl}^{\pi}M=\mathrm{Hom}\,(\Delta^1,\deleted{\unsure{\mathrm{Fl}}}M^{\pi})
 \qquad(\pi\geq\pi_0)
\]
forment aussi une filtration cardinale de $\underline{\mathrm{Fl}}\,M$. \deleted{\illegible{1}}

On dispose donc des notions de parties $\pi$-accessibles, accessibles, L-$\pi$-accessibles, L-accessibles de $\mathrm{Fl}\,M=\mathrm{Ob}\,\underline{\mathrm{Fl}}\,(M)$. \deleted{\unsure{Cela|Ces}} \unsure{les|Ces} notions sont stables par réunions et intersections finies, et les notions absolues (sans $\pi$) stables par intersections ``quelconques'' (i.e. \deleted{\unsure{pour une}} \inserted{pour} des \deleted{\unsure{cardin}} \uline{petites} familles de parties). Ce sera là \unsure{un des cas|une de mes} parti%
% archive sheet number: 200
\scanpage{202}%
\mspage{76}%
culièrement \unsure{crucials}, le cas d'une intersection $C\cap W$ !

Le résultat principal, pour lequel j'ai fait les \unsure{gammes} du § précédent, est le suivant :

\uline{Th. 5.1} Soit $M$ une catégorie satisfaisant les conditions dites :

% slanted
\margin{Démonstration \unsure{canulée|canular}, Énoncé peut-être \uuline{faux} !!} a) Stabilité par petites $\underrightarrow{\lim}$ (pas néc. filtr.

b) $M$ muni d'une filtration cardinale $(\mathrm{Filt}^{\unsure{\pi|R}})_{\pi\geq\pi_0}$,

c) les objets de $M$ sont accessibles

d) les $\mathrm{Filt}^\pi$ sont stables par $\underrightarrow{\lim}\!,$ indexées par des cat. $I$ telles que $\operatorname{Card}\mathrm{Fl}(I)\leq\pi$, et si $Z\in\mathrm{Filt}^\pi$, tout quotient $Z'$ de $Z$ (tel que $Z\longrightarrow Z'$ ne se factorise pas par un sous-objet $\unsure{Z'_0|Z'}$ de $Z'$ $\ne Z'$) est dans $\mathrm{Filt}^\pi$.

e) les limites inductives filtrantes dans $M$ transforment monos en monos.

\margin{$\big|$} Soit \deleted{\unsure{donc}} $C$ une partie L-accessible de $\mathrm{Fl}(M)=\mathrm{Ob}\,\underline{\mathrm{Fl}}\,(M)$, i.e. \inserted{\unsure{(5.17)|(5.11)}} une partie stable par $\underrightarrow{\lim}$ filtrantes, et accessible. \deleted{\unsure{Alors}} \inserted{Supposons de plus $C$, formé de monomorphismes\({}^{\ast}\), et \unsure{en} contenant les \uline{Iso}.}
% slanted
\margin{$\ast$ finalement, on a \unsure{eu besoin de|besoin de} cette hypothèse} Alors \deleted{soit} il existe une partie petite $C_0\subset C$, telle que $C\subset\widetilde C_0$. Plus précisément, si $\pi$ est un \deleted{\unsure{plus}} cardinal tel que $C$ soit $\pi$-accessible, et si on pose $C^\pi=C\cap\mathrm{Fl}^\pi M$, \deleted{de sorte qu'il} \inserted{alors on a}
% archive sheet number: 201
\scanpage{203}
\mspage{77}
\[
 C\subset\widetilde{C^{\unsure{\pi|R}}}=\widetilde C_0,
 \tag*{(5.1.1.)}
\]
où $C_0$ est n'importe quelle petite partie de $C^\pi$ telle que tout objet de $C^{\unsure{P|\pi}}$ soit isom. à un objet de $C_0$.

\uline{Cor. 6.2}. On a $C_{\unsure{\ast|\alpha}}=C_{0\ast}$, $\widetilde C=\widetilde C_0$.

\deleted{\uline{Cor. 6.3}. Pour qu'on ait $C=\widetilde C$, il f. et s. que $C$ soit stable par \unsure{cochg|cochgt}}

\uline{Cor. 6.3}. Conditions équivalentes sur \deleted{$C$ (\unsure{satisfais}} \inserted{la partie} L-accessible $C$ de $\mathrm{Fl}(M)_{\mathrm{mono}}$ \unsure{;} contenant les iso :

a) $C=\widetilde C$

b) $C$ est stable par cochgt. de base et \inserted{par composition \deleted{\unsure{de}}}.

b') $C$ est stable par cochgt de base, et par \unsure{sommes|somme} amalgam. \unsure{simple|simples} i.e. si $\deleted{\unsure{X\to}}\ A\unsure{-|\to} B_1$, $A-B_2$ sont dans $C$, alors $A-B_1\amalg_A B_2$ aussi.

c) $C_\pi$ est stable dans $\mathrm{Fl}^\pi M=\mathrm{Fl}\,(M^\pi)$ par cochgt de base et par composition

c') $C_\pi$ est stable dans $\mathrm{Fl}\,(M^\pi)$ par cochgt. de base et par sommes amalgamées (cf b').

Quand ces conditions sont satisfaites, alors $(C,C_\ast)$ est un couple de Quillen, justiciable du th. 1 du §1 (p. 10), \deleted{et de}

La dernière assertion résulte du th. 6.1. et du cor. au th. 1 du §1. On a
% archive sheet number: 202
\scanpage{204}
\mspage{78}
d'autre part tautologiquement
\diagram{\begin{tikzcd}[ampersand replacement=\&]
 a\arrow[r,Rightarrow]\&b\arrow[r,Rightarrow]\arrow[d,Rightarrow]\&b'\arrow[d,Rightarrow]\\
 \&c\arrow[r,Rightarrow]\&c'
\end{tikzcd}}
et il reste à prouver $c'\Longrightarrow a$, ce qu'on fera plus bas.

\unsure{Prouvons|Revenons} donc 6.1. Dans (6.1.1) l'égalité $\widetilde C_0=\widetilde{C^{\unsure{\pi|R}}}$ est tautologique. Donc tout revient à prouver
\[
 C\subset\widetilde{C^\pi}.
\]
une difficulté ici, c'est qu'on ne sait \uline{pas} que $\widetilde{C^\pi}$ soit stable par $\underrightarrow{\lim}$ filtrantes --- on le sait seulement pour $C$.

Soit donc
\[
 A\lhook\joinrel\longrightarrow B
\]
une flèche dans $C$, prouvons qu'elle est dans $\widetilde{C^\pi}$. Pour ceci, on \deleted{\unsure{pratique}} va construire, par induction transfinie, une suite transfinie \inserted{\uuline{\unsure{strictmt|strictemt}} croissante} de sous-objets $A_\alpha$ de $B$ contenant $A=A_0$, \deleted{et telle qu'on ait}
\diagram{\begin{tikzcd}[ampersand replacement=\&]
 A\arrow[r,hook,"{i_\alpha}"]\arrow[d,equal]\&A_\alpha\arrow[r,hook,"{\unsure{i_{\beta\alpha}|j_{\beta\alpha}}}"]\&\deleted{\unsure{B|A}}A_\beta\arrow[r,"{j_\beta}"]\&B\\
 A_0\&\&\&
\end{tikzcd}}
telle que l'on ait
\[
 i_\alpha\in\widetilde{C^{\pi}_{\unsure{4|\ast|\omega}}},\quad
 \inserted{\text{plus gén. les }i_{\beta,\alpha}\in\widetilde{C^\pi},}\quad j_\alpha\in C
\]
et que l'on ait
\[
 A_\alpha=\underset{\beta<\alpha}{\underrightarrow{\lim}}A_\beta
 \qquad\text{si }\deleted{\unsure{\beta}}\alpha\text{ ordinal limite.}
 \tag*{\unsure{(6.2.)|(6.2)}}
\]
% archive sheet number: 203
\scanpage{205}
\mspage{79}
Cela précise déjà la définition récurrente pour $\alpha\inserted{=\alpha_0}$ ordinal limite. Notons que le $A_{\alpha_0}$ construit en termes des $A_\beta$ $(\beta<\alpha_0)$ satisfait bien aux conditions que

a) les $\unsure{i_{\alpha_0,\beta}|i_{\beta,\beta}}:A_\beta\unsure{-|\to} A_{\alpha_0}$ sont dans $\widetilde C_{\unsure{\pi|R}}$ \deleted{\unsure{b)}} car $\widetilde C_\pi$ stable par ``limites inductives ordinales strictes'' (cf 1.19. p. 8).

b) $j_{\alpha_0}: A_{\alpha_0}\longrightarrow B$ est dans $C$, car c'est la limite inductive \inserted{filtrante} \deleted{\unsure{directe|d'un}} \unsure{du} système des $A_\beta\to B$ $(\beta<\alpha_0)$ qui sont dans $C$, et $C$ est stable par ces limites.

Il reste donc à préciser comment on passe de $A_\alpha$ à $A_{\alpha+1}$. Cela revient à prouver le

\uline{Lemme 6.4} Soit $f:A\lhook\joinrel\longrightarrow B$ dans $C$, \inserted{telle que $f$ ne soit pas un iso.} Alors il existe un sous-objet $A_1$ de $B$, \uline{distinct de} $A$, tel que $A\unsure{\to|-} A_1$ soit dans $\widetilde{C^\pi}$, et $A_1\to B$ dans $C$.

Si ce lemme est admis, cela signifie que la construction transfinie se continue aussi longtemps que l'on trouve des $A_\alpha$ tels que $A_\alpha\ne B$ i.e. $A_\alpha\hookrightarrow B$ non iso. Mais comme les $A_\alpha$ sont tous différents, si $\pi$ est l'ens des sous-objets de $B$ (on sait que c'est un
% archive sheet number: 204
\scanpage{206}
\mspage{80}
ens. petit, à cause des hyp. faites sur $M$), les ordinaux $\alpha$ pour lesquels $A_\alpha$ existe satisfont $\operatorname{card}A_\alpha\leq\unsure{\pi|R}$. Donc il doit exister un \deleted{$\unsure{A_{\alpha}}$} $\alpha$ tel que $A_\alpha=B$. Mais alors $i_\alpha\unsure{=|:|\doteq} A\longrightarrow A_\alpha$ est égal à $f$, donc $f\in\widetilde{C^\pi}$, qed

Il reste donc à prouver 6.4. Pour ceci, soit $I$ l'ens des sous-objets de $B$ qui sont dans $\mathrm{Filt}^\pi$ \deleted{$\unsure{=C^\pi}$}, ordonné par inclusion. (Si $B$ est dans $\mathrm{Filt}^\pi$, $A$ aussi, donc $f$ est \inserted{déjà} dans $C^\pi$\inserted{$\subset\unsure{C^\pi|\widetilde{C^\pi}}$}, et il n'y a rien à prouver.) C'est un ens. filtrant grand devant $\pi$, \inserted{(et $=$ $\pi$-\uline{adapté})} dont la $\underrightarrow{\lim}$ est $B$ (5.3. p. \unsure{45|46}). Considérons les
\[
 A_i=B_i\cap A
\]
ils sont dans $\mathrm{Filt}^\pi$, \deleted{\unsure{comme}} sous-objets d'objets de $\mathrm{Filt}^\pi$ (condition d) sur \unsure{la} filtration cardinale). Ainsi, la flèche $f$ apparaît comme limite inductive filtrante des flèches \deleted{$\unsure{A_i\to f_i}$}
\[
 f_i:A_i\to B_i .
\]
Il est \unsure{triviale|trivial} que le système des $A_i$ est cofinal dans $\mathrm{Ss}_\pi(A)$, donc
\[
 A=\underrightarrow{\lim}A_i,\quad\text{d'où}\quad f=\underrightarrow{\lim}f_i.
\]
% archive sheet number: 205
\scanpage{207}
\mspage{81}
\deleted{Mais} D'autre part, je dis que le syst. inductif des $f_i$ dans $\underline{\mathrm{Fl}}\,M$ est \uline{bien adapté à} $\unsure{\pi|R}$. On sait déjà que $I$ est \deleted{\unsure{a}} \unsure{bien}-adapté à $\pi$, le \deleted{\unsure{sup}} sup \inserted{dans $I$} d'une partie filtrante $J$ de cardinal $\leq\pi$ étant le sup dans l'ens. ordonné $\mathrm{Ss}_{\deleted{\unsure{\pi}}}(B)$ (lequel sup se trouve dans $\mathrm{Ss}_\pi(B)$). Il est essentiel, pour notre argument, de savoir que si $i_J=\sup J$, alors non seulement $B_{i_J}=\underset{i\in J}{\unsure{\underline{\lim}|\underrightarrow{\lim}}} B_i$, mais aussi $A_{i_J}=\underset{i\in J}{\underrightarrow{\lim}} A_i$, de sorte que \margin{$\Vert$}le système inductif des $f_i$ est $\pi$-adapté dans $\underline{\mathrm{Fl}}\,M$. Il semble qu'il faille donc une condition d'exactitude sur $M$ voisine de e, savoir

\textcircled{e'} les foncteurs $\underline{\lim}$ filtrantes commutent au chgt de base \deleted{\unsure{(dans}} \deleted{\unsure{(par un}}

(il \unsure{le|en} suffit ici dans le cas très particulier d'un chgt. de base par un mono \inserted{$A\to B$}, et d'une limite inductive filtr. de sous-objets de $B$ pour des morphismes \unsure{d'inclusions|d'inclusion} comme transitions. Notons que \textcircled{e}, \textcircled{e'} se trouvent généralement comme conséquences de la condition

(e$_0$) les $\underline{\lim}$ filtrantes dans $I$ sont exactes à gauche.
% archive sheet number: 206
\scanpage{208}
\mspage{82}
Ceci posé, on voit donc, $C$ étant $\unsure{\pi|R}$-accessible et $I$ grand devant $\pi$, que l'ens. $I_C$ des $i\in I$ tels que $f_i\in C$ est \uline{cofinal} dans $I$. Mais \deleted{les $f_i$} comme $f_i\in\mathrm{Filt}^\pi$, $f_i\in C$ \unsure{revient|équivaut} à $f_i\in C^\pi$. Comme $A\unsure{-|\longrightarrow}B$ n'est pas iso, il existe un $i_0$ tel que $B_{i_0}\not\subset A$, \deleted{\unsure{donc}} \inserted{puis} un $i\geq i_0$ tel que $f_i\in C^\pi$, \inserted{et aussi} $B_i\not\subset A$ (puisque $B_i\supset B_{i_0}$). Ainsi on a prouvé ceci

$(\ast)$ Sous les hyp. de (6.4.) et moyennant l'hyp. \textcircled{e'} sur $M$, il existe un sous-objet $B'$ de $B$ satisfaisant les conditions suivantes

a) $B'\in\mathrm{Filt}^\pi$

b) $B'\not\subset A$

c) Posant $A'=B'\cap A$, l'inclusion
\[
 i':A'-B'
\]
est dans $C^\pi$.

De plus, l'ens des $J\subset\mathrm{Ss}_\pi(B)$ formé des $B'$ en question est cofinal dans $\mathrm{Ss}_\pi(B)$, et $=$ $\pi$-cofinal.

Il faut prouver encore ce dernier point --- que $J$ (\deleted{\unsure{dont}} on sait qu'il est cofinal) est $=$ $\pi$-cofinal, \deleted{\unsure{il est donc}} i.e. stable par $\sup$ \inserted{(dans $\mathrm{Ss}_\pi(B)$)} de \deleted{\unsure{ses}} parties filtrantes $J'\in J$ telles que $\operatorname{card}J'\leq\pi$. Mais cela résulte aussitôt du fait que $C$ étant stable par
% archive sheet number: 207
\scanpage{209}
\mspage{83}
limites inductives filtrantes, $C^{\unsure{\pi|R}}=C\cap\mathrm{Filt}^\pi$ est stable par $\underrightarrow{\lim}$ filtrantes de cardinal $\leq\pi$, en vertu de la condition d) sur $\mathrm{Filt}^\pi$.

Considérons maintenant le carré cocartésien
\diagram{\begin{tikzcd}[ampersand replacement=\&]
 A'\arrow[r,hook',"i'"]\arrow[d,hook]\&B'\arrow[d]\&\\
 A\arrow[r,hook',"{\deleted{\unsure{g}}i}"]\& C'=A\amalg_{A'}B'\arrow[r,"{\unsure{j|J}}"]\&B
\end{tikzcd}}
Comme $i'\in C^\pi$, on a $i\in\widetilde{C^\pi}$ (car $\widetilde{C^\pi}$ stable par cochgt de base). D'autre part, j'ai besoin du fait que $j$ soit dans $C$, \deleted{\unsure{a priori} \unsure{plu}} et je \uline{ne} sais pas $=$ qu'il est mono ! \unsure{Aussi} il faut repasser aux notations indicielles
\diagram{\begin{tikzcd}[ampersand replacement=\&]
 A_i\arrow[r,hook,no head,"\alpha_i"]\arrow[d,hook,no head]\&B_i\arrow[d]\&\\
 A\arrow[r,hook,"\beta_i"]\&C_i\arrow[r,"{\unsure{\gamma_i|\delta_i}}"]\&B
\end{tikzcd}}
Comme \unsure{le foncteur|la formation} des sommes amalgamées commute aux limites inductives filtrantes, on sait que
\[
 \begin{aligned}
 \unsure{\underline{\lim}|\underrightarrow{\lim}}\,^{\unsure{\prime\prime}} C_i
 &=\deleted{\underrightarrow{\lim}A_i}\ A\amalg_{\underrightarrow{\lim} A_i}\underline{\lim}B_i\\
 &=A\amalg_A B=B,
 \end{aligned}
\]
donc la limite inductive des $\gamma_i$ est un iso. (\unsure{can.} $\mathrm{id}_B$). \unsure{D'autre part}, le $=$ argument montre que le système inductif des $C_i$ $(i\in J)$ \inserted{de $M$} est \deleted{lui aussi} $\pi$-adapté, donc aussi celui des $\gamma_i$.$^{\ast}$

% slanted
\margin{\uuline{Sauf} qu'on ne sait pas que les $C_i$ soient dans $\mathrm{Filt}^\pi$ --- en fait ils ne seront \unsure{guère|pas} en général, \illegible{2} \unsure{seraient|serait} \unsure{supposer|supposés}, \unsure{puisque} $A,B\in\mathrm{Filt}^{\unsure{\pi}}$}
% archive sheet number: 208
\scanpage{210}
\mspage{84}
dans $\underline{\mathrm{Fl}}(M)$. Comme $\underrightarrow{\lim}\gamma_i$ \deleted{$\unsure{L|i\in J}$} \deleted{$\unsure{=B}$} est un iso, il est dans $C$. \deleted{Car la stabilité de $C$ implique que $C$} Comme $C$ est $\unsure{\pi|R}$-accessible, il s'ensuit que l'ens $J'$ formé des \deleted{$\unsure{J}$} $i\in J$ tels que $\gamma_i\in C$ est lui aussi cofinal.$^{+}$
% slanted
\margin{\uuline{canulé}, cf annotation précédente.} Donc on trouve $i$ avec $\alpha_i\in C^\pi$, $\gamma_i\in C$, ce qui établit le lemme 6.4, et achève la démonstration du théorème 6.1,

\deleted{Chemin faisant j'ai de plus en plus l'impression que le rôle que je fais jouer dans \unsure{tout ça|tout §ci} aux monomorphismes, et surtout l'hypothèse que $C$ soit formé de monomorphismes, est tout à fait artificiel. Je vais essayer de m'en débarrasser.}

7. \uline{Construction de triples de Quillen clos avec} $C=\mathrm{Mon}\subset\mathrm{Cof}_W$

Finalement, tous les beaux développements des §§ 5,6 ne me donnent pas le ``sweeping result'' souhaité. Peut-être est-il faux --- peut-être la propriété d'engendrement $\sim$ par une petite partie (pour les ensembles accessibles, ou L-accessibles) est-elle beaucoup plus \unsure{rarement|souvent} satisfaite, que je ne l'aurais imaginé. Pour bien faire, il faudrait que je traite jusqu'au
% archive sheet number: 209
\scanpage{211}
\mspage{85}
bout quelques exemples, p.ex. celui de $\mathrm{Cof}_{W}$ et de $W$-même dans $\mathrm{Cat}$, pour un localiseur fondamental plus ou moins général, et surtout dans le cas $W=W_\infty$. Sûrement $\mathrm{Cof}_{W}$ et $W$-même sont accessibles (il faudrait que je le vérifie pourtant), et si le ``théorème'' hypothétique du \S~6 est faux pour eux, au moins en aurai-je le cœur net. Mais il reste pourtant la question : si $C=\widetilde C$ et stable par $\underrightarrow{\lim}$ filtrantes et de la forme $\widetilde C_0$, $C_0$ petit, et si $W$ est $\unsure{\underline L|L}$-accessible également, alors $C\cap W$ est-il \deleted{aussi} de la forme $\widetilde{\unsure{TC|T|\pi}}_0$, avec $TC_0$ petit ? Je viens \unsure{encore|en vain} d'essayer d'y \unsure{penser|pousser|passer}, en supposant de plus $C\subset\mathrm{monos}$, et ne vois aucun moyen d'y parvenir. Il faudrait donc un contreexemple pour dissiper ces perplexités ! Et le mieux serait dans $\mathrm{Cat}$, tant qu'à faire.

Il surnage pourtant un résultat \inserted{tangible} positif, que je vais maintenant essayer de bien cerner et prouver :
% archive sheet number: 210
